\documentclass{acm_proc_article-sp} \usepackage{graphicx} \usepackage{amsmath} \usepackage{booktabs} \usepackage{url}

\begin{document} \title{AI201 PA3 Report\ Implementing the Backpropagation Algorithm:\ Multilayer Perceptron for 8-Class Classification}

\numberofauthors{1} \author{ \alignauthor Marcus Rafael B. Tiongson (2013-59474)\ ~\ Date of Submission: March 27, 2026 }

\maketitle

% ----------------------------------------------------------------------- \section{Introduction} % -----------------------------------------------------------------------

This report documents the implementation and evaluation of a Multilayer Perceptron (MLP) neural network trained via the Backpropagation Algorithm on an 8-class dataset. The implementation uses only the \texttt{numpy} library for all forward and backward propagation computations, as required, with \texttt{matplotlib} and \texttt{csv} used for visualization and file I/O respectively.

The dataset consists of 3,486 instances with 354 numerical features per instance. The class distribution is highly imbalanced: Class~1 accounts for roughly 46.6% of all samples, while Class~4 constitutes the smallest share. Without intervention, such skew causes a network to collapse toward a trivial majority-class predictor. The Synthetic Minority Over-sampling Technique (SMOTE)~\cite{chawla2002smote} is therefore applied prior to training to construct a balanced dataset.

Two network configurations are explored under a shared architecture of two hidden layers and one logistic output layer. Network~A uses the hyperbolic tangent (Tanh) activation function for both hidden layers; Network~B uses the Leaky Rectified Linear Unit (Leaky ReLU) for both hidden layers. Both networks use the logistic sigmoid as the output activation. The output layer has 8 neurons, one per class, and targets are one-hot encoded.

Training follows the Generalized Delta Rule with momentum, applied in mini-batches of size 8. Ten configurations in total are evaluated---five per network family---through a systematic sweep over the learning rate~$\eta$, momentum coefficient~$\alpha$, and hidden layer size. Each configuration is evaluated on a fixed 800-sample validation set using macro-averaged F1 score, Matthews Correlation Coefficient (MCC), and time-adjusted training cost.

% ----------------------------------------------------------------------- \section{Objectives} % -----------------------------------------------------------------------

This work aims to implement the backpropagation algorithm entirely from scratch in NumPy, supporting user-configurable activation functions per layer and mini-batch gradient descent with a default batch size of 8. The implementation further addresses the severe class imbalance in the training data through SMOTE augmentation applied before partitioning.

The primary experimental objective is to systematically tune the learning rate~$\eta$, momentum coefficient~$\alpha$, and hidden layer size for both Network~A (Tanh) and Network~B (Leaky ReLU), and to identify the configuration that achieves the highest macro-averaged F1 score at the lowest time-adjusted training cost. A secondary objective is to compare the convergence behavior and classification performance of the Tanh and Leaky ReLU activation families on this dataset, and to generate labeled predictions for the held-out unlabeled test set using the best trained models from each family.

% ----------------------------------------------------------------------- \section{Methodology} % -----------------------------------------------------------------------

\subsection{Dataset and Class Balancing}

The dataset comprises 3,486 instances with 354 continuous features and class labels ranging from 1 to 8. As shown in Figure~\ref{fig:dist_original}, Class~1 dominates the distribution at 46.6% of all samples, introducing a substantial risk of majority-class bias. SMOTE~\cite{chawla2002smote} is applied to the full dataset (random state 50) to generate synthetic minority-class samples via linear interpolation between existing instances in feature space, yielding a fully balanced distribution across all eight classes (Figure~\ref{fig:dist_smote}). One limitation of this approach is that SMOTE is applied \emph{before} the train-validation split, meaning some synthetic samples may appear in the validation set. This introduces a mild optimistic bias into validation metrics and would ideally be corrected by applying SMOTE strictly within the training fold after the split.

Following balancing, the dataset is partitioned into a training set and a fixed validation set of exactly 800 samples (random state 50). The remaining samples form the training set. All four partitions are saved to CSV files for reproducibility.

\begin{figure}[ht] \centering \includegraphics[width=0.85\columnwidth]{figures/Class_Distribution_Original_Dataset.png} \caption{Class distribution of the original dataset. Class~1 accounts for 46.6% of all instances.} \label{fig:dist_original} \end{figure}

\begin{figure}[ht] \centering \includegraphics[width=0.85\columnwidth]{figures/Class_Distribution_after_SMOTE.png} \caption{Class distribution after SMOTE augmentation. Each class contains an equal number of samples.} \label{fig:dist_smote} \end{figure}

\subsection{Network Architecture}

The MLP consists of two hidden layers (layer~$i$ and layer~$j$) followed by an output layer (layer~$k$). The input to the first hidden layer is the raw feature vector augmented with a prepended bias unit of value $+1$, giving an effective input dimensionality of 355. Each hidden layer prepends its own bias to its output before passing it to the next layer. The output layer has 8 neurons, one per class. All weight matrices include the bias weight in column~0, so no separate bias vector is required.

Weight matrices are initialized by sampling from a zero-mean Gaussian distribution with standard deviation~0.01, using random seed~50 for reproducibility: [ w^{(l)}_{ji}(0) \sim \mathcal{N}(0,;0.01^2). ]

\subsection{Forward Propagation}

For each layer $l \in {1, \dots, L}$, the internal activity (pre-activation) of neuron~$j$ is the weighted sum of inputs from the previous layer, where index $i=0$ carries the bias weight $w^{(l)}_{j0}$ with fixed input $\varphi_0 = +1$:

\begin{equation} v^{(l)}*j(n) = \sum*{i=0}^{p} w^{(l)}_{ji}(n),\varphi_i!\left(v^{(l-1)}_i(n)\right) \label{eq:internal_activity} \end{equation}

The neuron output is obtained by applying the activation function $\varphi_j$:

\begin{equation} y_j(n) = \varphi_j!\left(v^{(l)}_j(n)\right) \label{eq:neuron_output} \end{equation}

\subsection{Activation Functions}

Derivatives are expressed directly in terms of the neuron output $o(n) = \varphi(v(n))$ to avoid recomputing the pre-activation.

\textbf{Logistic sigmoid} (output layer, $a = 1.0$): \begin{equation} \varphi(v) = \frac{1}{1 + e^{-av}}, \qquad \varphi'(v) = a \cdot o(n),\bigl(1 - o(n)\bigr) \label{eq:logistic} \end{equation}

\textbf{Hyperbolic tangent} (Network~A hidden layers, $a = 1.716$, $b = 2/3$): \begin{equation} \varphi(v) = a\tanh(bv), \qquad \varphi'(v) = \frac{b}{a}\bigl(a - o(n)\bigr)!\bigl(a + o(n)\bigr) \label{eq:tanh} \end{equation}

The parameter values $a = 1.716$ and $b = 2/3$ are recommended in the literature~\cite{lecun1998efficient} to produce outputs with unit variance and decorrelated gradients, which accelerates convergence.

\textbf{Leaky ReLU} (Network~B hidden layers, $\gamma = 0.01$): \begin{equation} \varphi(v) = \begin{cases} v & v > 0 \ \gamma,v & v \leq 0 \end{cases}, \qquad \varphi'(v) = \begin{cases} 1 & o(n) > 0 \ \gamma & o(n) \leq 0 \end{cases} \label{eq:lrelu} \end{equation}

The small negative-side slope $\gamma = 0.01$ prevents dying neurons while keeping the activation nearly linear.

\subsection{Error Signal and Mean Square Error}

The error at output neuron~$k$ is the difference between the desired one-hot target $d_k$ and the network output $o_k(n) = \varphi_k(v^{(L)}_k(n))$:

\begin{equation} e_k(n) = d_k - o_k(n) \label{eq:error_signal} \end{equation}

The instantaneous mean square error (MSE) over all output neurons $\mathcal{C}$ is:

\begin{equation} \mathcal{E}(n) = \frac{1}{2} \sum_{k \in \mathcal{C}} e_k^2(n) \label{eq:mse} \end{equation}

\subsection{Backpropagation: Delta Computation}

For the output layer~$L$, the local gradient (delta) is:

\begin{equation} \delta^{(L)}_k(n) = e_k(n) \cdot \varphi'!\left(v^{(L)}_k(n)\right) \label{eq:delta_output} \end{equation}

For each hidden layer $l \in {L-1, L-2, \dots, 1}$, the delta is computed by propagating the error backward:

\begin{equation} \delta^{(l)}_j(n) = \varphi'!\left(v^{(l)}*j(n)\right) \cdot \sum*{k}, \delta^{(l+1)}*k(n), w^{(l+1)}*{kj}(n) \label{eq:delta_hidden} \end{equation}

The bias column (column~0) of the weight matrix is excluded from the backward sum, since it does not correspond to a neuron in the previous layer.

\subsection{Weight Update: Generalized Delta Rule with Momentum}

Gradients are accumulated over each mini-batch of $B$ samples before the weight update is applied. The weight increment at step~$n$ follows the Generalized Delta Rule with a momentum term~$\alpha$ that adds a fraction of the previous increment:

\begin{equation} \Delta w^{(l)}*{ji}(n) = \alpha,\Delta w^{(l)}*{ji}(n-1) ;+; \eta,\delta^{(l)}_j(n),\varphi_i!\left(v^{(l-1)}_i(n)\right) \label{eq:weight_increment} \end{equation}

\begin{equation} w^{(l)}*{ji}(n+1) = w^{(l)}*{ji}(n) + \Delta w^{(l)}_{ji}(n) \label{eq:weight_update} \end{equation}

Bias weights ($i = 0$) are updated identically, with $\varphi_0 = +1$. When $\alpha = 0$, the rule reduces to the standard gradient descent update. For the Leaky ReLU layers, gradient norms are clipped at 5.0 to prevent overshooting in early training epochs.

\subsection{Mini-Batch Training and Evaluation}

Training iterates over the training set in sequential mini-batches of $B = 8$ samples per batch. Within each batch, forward and backward passes are performed for every sample and the gradients are accumulated; the accumulated gradient is divided by $B$ before the weight update. Training and validation MSE are logged at every epoch. Validation MSE and classification scores (accuracy, per-class precision, recall, F1, and MCC) are computed on the fixed 800-sample validation set at the end of training.

\subsection{Hyperparameter Configurations}

Table~\ref{tab:configs} summarizes the ten configurations evaluated. Each run trains for 100 epochs. The column ``Epoch (MSE $<$ 0.02)'' records the first epoch at which the training MSE falls below the 0.02 target, used subsequently to compute the time-adjusted training cost.

\begin{table}[ht] \centering \caption{Hyperparameter configurations for Networks A and B.} \label{tab:configs} \small \begin{tabular}{llccc} \toprule Network & Config & $\eta$ & $\alpha$ & Size \ \midrule A (Tanh) & Baseline & 0.85 & 0.9 & 8 \ A (Tanh) & improv & 0.10 & 0.5 & 12 \ A (Tanh) & improv2 & 0.05 & 0.5 & 12 \ A (Tanh) & improv2_fast & 0.05 & 0.9 & 12 \ A (Tanh) & improv2_small & 0.05 & 0.9 & 5 \ \midrule B (Leaky ReLU) & Baseline & 0.85 & 0.9 & 8 \ B (Leaky ReLU) & improv & 0.10 & 0.5 & 12 \ B (Leaky ReLU) & improv2 & 0.05 & 0.5 & 12 \ B (Leaky ReLU) & improv2_fast & 0.05 & 0.9 & 12 \ B (Leaky ReLU) & improv2_small & 0.05 & 0.9 & 5 \ \bottomrule \end{tabular} \end{table}

% ----------------------------------------------------------------------- \section{Experimental Results} % -----------------------------------------------------------------------

\subsection{Network A (Tanh)}

Training begins with the default configuration ($\eta = 0.85$, $\alpha = 0.9$, hidden size 8). The training MSE remains essentially flat across all 100 epochs with no discernible descent, and the resulting macro-averaged F1 score is 0.2261. Inspection of the confusion matrix confirms that the network collapses to predicting a single class exclusively, consistent with weight updates that overshoot the loss surface due to the excessively large learning rate. The learning curve for this baseline run is shown in Figure~\ref{fig:lcA_baseline}.

\begin{figure}[ht] \centering \includegraphics[width=0.95\columnwidth]{figures/Learning_curve_Network_A.png} \caption{Learning curve for Network~A baseline ($\eta=0.85$, $\alpha=0.9$, size=8). Training error remains flat, indicating no learning.} \label{fig:lcA_baseline} \end{figure}

Configuration A-improv ($\eta = 0.1$, $\alpha = 0.5$, size = 12) reduces the learning rate to 0.1 and the momentum to 0.5, while increasing the hidden layer size to 12. This is sufficient to achieve convergence: both training and validation MSE descend steadily, though the validation curve exhibits noticeable oscillations. The training MSE crosses below the 0.02 target at approximately epoch 72, with a macro-averaged F1 score of 0.9730. The learning curve is shown in Figure~\ref{fig:lcA_improv}.

\begin{figure}[ht] \centering \includegraphics[width=0.95\columnwidth]{figures/Learning_curve_Network_A_improved.png} \caption{Learning curve for Network~A improved ($\eta=0.1$, $\alpha=0.5$, size=12). Convergence is achieved; the validation error is more erratic than the training error.} \label{fig:lcA_improv} \end{figure}

Configuration A-improv2 ($\eta = 0.05$, $\alpha = 0.5$, size = 12) halves the learning rate again, yielding the highest final F1 of all ten configurations at 0.9860, but requires the full 100 epochs to converge (epoch 97). The validation error is more erratic in this run than in A-improv, suggesting diminishing returns in convergence smoothness at this momentum level.

Configuration A-improv2_fast ($\eta = 0.05$, $\alpha = 0.9$, size = 12) restores momentum to 0.9 while retaining $\eta = 0.05$. This dramatically accelerates convergence: the training MSE falls below 0.02 at epoch 35. The final macro-averaged F1 is 0.9853, marginally lower than A-improv2 but nearly identical in practice. Applying early stopping at epoch 35 yields an adjusted training cost of approximately 24 seconds. The learning curve is shown in Figure~\ref{fig:lcA_fast}.

\begin{figure}[ht] \centering \includegraphics[width=0.95\columnwidth]{figures/Learning_curve_Network_A_improved_(Fast).png} \caption{Learning curve for Network~A improved fast ($\eta=0.05$, $\alpha=0.9$, size=12). Convergence below MSE~0.02 occurs at epoch~35.} \label{fig:lcA_fast} \end{figure}

Configuration A-improv2_small ($\eta = 0.05$, $\alpha = 0.9$, size = 5) reduces the hidden layer size to 5 nodes with the same favorable learning rate and momentum settings. Convergence slows considerably: the training MSE does not reliably fall below 0.02 within 100 epochs, and the validation error is noticeably more erratic, yielding a macro-averaged F1 of 0.9616. This is consistent with insufficient representational capacity for the 354-dimensional input space.

\subsection{Network B (Leaky ReLU)}

The baseline configuration for Network~B ($\eta = 0.85$, $\alpha = 0.9$, size = 8) produces a flat MSE curve across all 100 epochs and an F1 score of 0.0250, indicating complete failure to learn. The behavior mirrors Network~A's baseline. The learning curve is shown in Figure~\ref{fig:lcB_baseline}.

\begin{figure}[ht] \centering \includegraphics[width=0.95\columnwidth]{figures/Learning_Curve_Network_B.png} \caption{Learning curve for Network~B baseline ($\eta=0.85$, $\alpha=0.9$, size=8). No convergence; error remains at its initial value throughout.} \label{fig:lcB_baseline} \end{figure}

Configuration B-improv ($\eta = 0.1$, $\alpha = 0.5$, size = 12) restores convergence, with training MSE crossing the 0.02 threshold at approximately epoch 76 and a macro-averaged F1 of 0.9656. However, the validation error diverges from the training error toward the end of the run, driven largely by a high number of false positives for Class~6 (19 false positives in the validation confusion matrix). The learning curve is shown in Figure~\ref{fig:lcB_improv}.

\begin{figure}[ht] \centering \includegraphics[width=0.95\columnwidth]{figures/Learning_curve_Network_B_improved.png} \caption{Learning curve for Network~B improved ($\eta=0.1$, $\alpha=0.5$, size=12). Convergence is achieved but the validation error diverges from the training error toward epoch 100.} \label{fig:lcB_improv} \end{figure}

Configuration B-improv2 ($\eta = 0.05$, $\alpha = 0.5$, size = 12) reduces the learning rate to 0.05, improving the macro F1 to 0.9778 and partially stabilizing the validation curve. Weight improvement begins around epoch 20. Classes~6 and~8 continue to contribute the most classification errors.

Configuration B-improv2_fast ($\eta = 0.05$, $\alpha = 0.9$, size = 12) restores $\alpha = 0.9$, which reduces Class~6 false positives from 19 to 3. The training error approaches the 0.02 target at approximately epoch 65, yielding a macro-averaged F1 of 0.9819 and an adjusted training time of approximately 48 seconds. The validation error is noticeably more stable than the training error in this run, suggesting reasonable generalization.

Configuration B-improv2_small ($\eta = 0.05$, $\alpha = 0.9$, size = 5) produces a macro F1 of 0.9143 and MCC of 0.9018, confirming that representational capacity is a limiting factor for Network~B as well.

\subsection{Consolidated Learning Curves}

Figure~\ref{fig:lc_consolidated} presents all converging configurations overlaid for direct comparison.

\begin{figure}[ht] \centering \includegraphics[width=0.95\columnwidth]{figures/Consolidated_Learning_Curves.png} \caption{Consolidated learning curves for all converging configurations across Networks~A and~B.} \label{fig:lc_consolidated} \end{figure}

\subsection{Summary of Results}

Table~\ref{tab:results_f1} presents all ten configurations ranked by macro-averaged F1 score. Table~\ref{tab:results_time} presents the eight converging configurations ranked by time-adjusted training cost, computed as:

\begin{equation} T_{\text{adj}} = T_{\text{total}} \times \frac{n_{\text{converge}}}{100} \label{eq:time_adj} \end{equation}

where $n_{\text{converge}}$ is the epoch at which training MSE first falls below 0.02. Configurations that never reach this threshold (DNF) receive no adjusted time. Configurations requiring beyond 100 epochs use a multiplier of 1.0.

\begin{table}[ht] \centering \caption{All configurations ranked by macro-averaged F1 score.} \label{tab:results_f1} \small \begin{tabular}{clccccc} \toprule Rank & Network & Epoch & Acc. & F1 & MCC & Time (s) \ \midrule 1 & A-improv2 & 97 & 0.9863 & 0.9860 & 0.9843 & 62.36 \ 2 & A-improv2_fast & 35 & 0.9850 & 0.9853 & 0.9829 & 68.34 \ 3 & B-improv2_fast & 65 & 0.9813 & 0.9819 & 0.9787 & 74.51 \ 4 & B-improv2 & 97 & 0.9775 & 0.9778 & 0.9745 & 76.56 \ 5 & A-improv & 72 & 0.9725 & 0.9730 & 0.9689 & 65.33 \ 6 & B-improv & 76 & 0.9650 & 0.9656 & 0.9603 & 73.22 \ 7 & A-improv2_small & 100+ & 0.9613 & 0.9616 & 0.9561 & 63.41 \ 8 & B-improv2_small & 100+ & 0.9138 & 0.9143 & 0.9018 & 71.15 \ 9 & A-baseline & DNF & 0.1275 & 0.2262 & --- & 63.60 \ 10 & B-baseline & DNF & 0.1113 & 0.0250 & 0.000 & 71.80 \ \bottomrule \end{tabular} \end{table}

\begin{table}[ht] \centering \caption{Converging configurations ranked by time-adjusted training cost.} \label{tab:results_time} \small \begin{tabular}{clcccc} \toprule Rank & Network & Epoch & Time (s) & $T_{\text{adj}}$ (s) & Config \ \midrule 1 & A-improv2_fast & 35 & 68.34 & \textbf{23.92} & $\eta$=0.05, $\alpha$=0.9, sz=12 \ 2 & A-improv & 72 & 65.33 & 47.04 & $\eta$=0.1, $\alpha$=0.5, sz=12 \ 3 & B-improv2_fast & 65 & 74.51 & 48.43 & $\eta$=0.05, $\alpha$=0.9, sz=12 \ 4 & B-improv & 76 & 73.22 & 55.65 & $\eta$=0.1, $\alpha$=0.5, sz=12 \ 5 & A-improv2 & 97 & 62.36 & 60.49 & $\eta$=0.05, $\alpha$=0.5, sz=12 \ 6 & A-improv2_small & 100+ & 63.41 & 63.41 & $\eta$=0.05, $\alpha$=0.9, sz=5 \ 7 & B-improv2 & 97 & 76.56 & 74.27 & $\eta$=0.05, $\alpha$=0.5, sz=12 \ 8 & B-improv2_small & 100+ & 71.15 & 71.15 & $\eta$=0.05, $\alpha$=0.9, sz=5 \ \bottomrule \end{tabular} \end{table}

% ----------------------------------------------------------------------- \section{Analysis and Discussion of Results} % -----------------------------------------------------------------------

\subsection{Effect of Learning Rate}

Reducing $\eta$ from 0.85 to 0.1 is the single most critical intervention in this study. Both network families fail completely at the default learning rate: weight updates are so large that the loss surface is perpetually overshot, and the network collapses to a single-class prediction within the first few epochs. At $\eta = 0.1$, convergence is restored for both families. A further reduction to 0.05 yields marginally higher F1 scores---by approximately 0.01 to 0.013 for Tanh configurations---but requires additional epochs to reach the same training MSE threshold. The relationship is consistent with the standard tradeoff for step size in gradient descent: smaller steps reduce oscillation and allow finer-grained minimization but come at the cost of slower progress per epoch.

\subsection{Effect of Momentum}

Momentum plays a dual role. During the initial phase of tuning, $\alpha$ is reduced from 0.9 to 0.5 alongside the learning rate reduction, ensuring that the first converging configurations do not inherit the instability of the default setup. Once a stable learning rate of 0.05 is established, restoring $\alpha$ to 0.9 provides a large convergence speedup: for Network~A, the epoch of convergence decreases from 97 to 35 (a 64% reduction); for Network~B, from approximately 97 to 65 (33% reduction). Crucially, this speed gain comes at no cost to final accuracy: F1 scores for the \emph{fast} variants are within 0.001 of their slower counterparts.

\subsection{Effect of Hidden Layer Size}

Compact networks with hidden size 5 consistently perform worse than their size-12 counterparts across both activation families. For Network~A, the F1 drops from 0.9853 to 0.9616; for Network~B, from 0.9819 to 0.9143. The degradation is more severe for Network~B, suggesting that Leaky ReLU requires greater representational width to partition the 354-dimensional feature space reliably. In both cases, the convergence epoch exceeds 100, confirming that size-5 hidden layers provide insufficient capacity for this task.

\subsection{Tanh versus Leaky ReLU}

Tanh (Network~A) marginally outperforms Leaky ReLU (Network~B) across all matched hyperparameter configurations. The margin is approximately 0.003 to 0.007 F1 points for the well-tuned configurations. The most notable difference is in Class~6 classification: Network~B configurations exhibit a systematically higher false positive rate for Class~6, most pronounced at 19~FP for B-improv and reduced to 3~FP for B-improv2_fast after momentum restoration. The Tanh activation's bounded, zero-centered outputs are known to produce better-conditioned gradients than the unbounded outputs of Leaky ReLU~\cite{lecun1998efficient}, which may explain the more stable class discrimination across all eight classes. Network~B is consistently approximately 8 to 10 seconds slower per 100 epochs than Network~A with equivalent hyperparameters, attributable to the per-layer gradient norm clipping applied to Leaky ReLU weight updates.

\subsection{Model Quality and Confusion Matrix}

The recommended model, A-improv2_fast, achieves a validation accuracy of 0.9850 with an MCC of 0.9829. The MCC near 1.0 confirms that the model performs well across all class pairs without systematic bias. Per-class F1 scores are uniformly high across all 8 classes, with all individual class values exceeding 0.96. The backup model, B-improv2_fast, achieves a validation accuracy of 0.9813 and MCC of 0.9787, with similarly uniform per-class performance except for the residual Class~6 sensitivity already noted.

\subsection{Data Leakage Caveat}

The validation metrics reported throughout this study carry mild optimistic bias because SMOTE was applied before the train-validation split. Synthetic samples generated from minority-class training instances may appear in the validation set, potentially inflating accuracy and F1 scores for those classes. In a rigorously designed evaluation pipeline, SMOTE would be applied exclusively within the training fold after each split. This limitation is acknowledged but accepted for simplicity in the current implementation.

% ----------------------------------------------------------------------- \section{Conclusion} % -----------------------------------------------------------------------

This study implemented and evaluated ten configurations of a two-hidden-layer MLP trained via backpropagation with momentum on an 8-class imbalanced dataset, using two activation function families (Tanh and Leaky ReLU) and systematic variation of $\eta$, $\alpha$, and hidden layer size. Four principal findings emerge.

First, the default hyperparameters ($\eta = 0.85$, $\alpha = 0.9$) are categorically unsuitable for both network families: neither achieves any meaningful gradient descent, yielding F1 scores of 0.2261 and 0.0250. Reducing $\eta$ to 0.1 is the necessary first intervention. Second, once a stable learning rate is established, restoring $\alpha = 0.9$ provides a large convergence speedup with no accuracy penalty; for Network~A, the convergence epoch decreases from 97 to 35. Third, a hidden layer size of 12 is required for reliable performance on the 354-dimensional input; size-5 networks fail to converge within 100 epochs and exhibit noticeably lower F1 scores. Fourth, Tanh marginally outperforms Leaky ReLU on this dataset, with consistently superior F1 scores and fewer class-specific misclassifications.

The recommended model for deployment is \textbf{A-improv2_fast} (Tanh, $\eta = 0.05$, $\alpha = 0.9$, size = 12), which achieves a macro-averaged F1 of 0.9853 and an adjusted training cost of approximately 24 seconds with early stopping at epoch 35. Its trained weights are stored in \texttt{modelA/trained_weights.csv}. The backup model is \textbf{B-improv2_fast} (Leaky ReLU, $\eta = 0.05$, $\alpha = 0.9$, size = 12) with a macro-averaged F1 of 0.9819 and an adjusted training cost of approximately 48 seconds. For future work, SMOTE should be applied within the training fold after the split to eliminate data leakage, and formal early stopping based on validation error should be implemented to reduce unnecessary computation.

% ----------------------------------------------------------------------- \section{References} % -----------------------------------------------------------------------

\begin{thebibliography}{9}

\bibitem{haykin2009} S.~Haykin, \newblock \emph{Neural Networks and Learning Machines}, 3rd ed. \newblock Prentice Hall, 2009.

\bibitem{chawla2002smote} N.~V.~Chawla, K.~W.~Bowyer, L.~O.~Hall, and W.~P.~Kegelmeyer, \newblock SMOTE: Synthetic minority over-sampling technique. \newblock \emph{Journal of Artificial Intelligence Research}, 16:321--357, 2002.

\bibitem{lecun1998efficient} Y.~LeCun, L.~Bottou, G.~B.~Orr, and K.-R.~M"{u}ller, \newblock Efficient backprop. \newblock In \emph{Neural Networks: Tricks of the Trade}, Springer, 1998.

\end{thebibliography}

\end{document}
